\documentclass[a4paper,12px]{article}

\usepackage{graphicx}
\usepackage[english]{babel}
\usepackage{fancyhdr}
\usepackage{lastpage}
\usepackage{xifthen}
\usepackage[linesnumberedhidden, titlenotnumbered]{algorithm2e}
\usepackage{lipsum}
\usepackage{hyperref}
\usepackage{array}
\usepackage{tabularx}

\usepackage{minted}
\usepackage{caption}
\usepackage{amssymb}

\pagestyle{fancy}
\lhead{\includegraphics[width=7cm]{logoUvA}}
\rhead{\footnotesize \textsc {Paper\\ \titel}}
\lfoot
{
    \footnotesize \studentA
    \ifthenelse{\isundefined{\studentB}}{}{\\ \studentB}
    \ifthenelse{\isundefined{\studentC}}{}{\\ \studentC}
    \ifthenelse{\isundefined{\studentD}}{}{\\ \studentD}
    \ifthenelse{\isundefined{\studentE}}{}{\\ \studentE}
}
\cfoot{}
\rfoot{\small \textsc {Page \thepage\ of \pageref{LastPage}}}
\renewcommand{\footrulewidth}{0.5pt}

\fancypagestyle{firststyle}
{
    \vspace*{6cm}
    \fancyhf{}
    \renewcommand{\headrulewidth}{0pt}
    \chead{\includegraphics[width=7cm]{logoUvA}}
    %\rfoot{\small \textsc {Page \thepage\ of \pageref{LastPage}}}
}

\setlength{\topmargin}{-0.3in}
\setlength{\textheight}{630pt}
\setlength{\headsep}{40pt}

% =================================== DOC INFO ===================================

\newcommand{\titel}{Trade-offs in distributed file systems}
\newcommand{\opdracht}{Examples and applications}
\newcommand{\docent}{Ana Varbanescu}
\newcommand{\cursus}{Concurrency and Parallel Programming}
\newcommand{\vakcode}{5062COPP6Y}
\newcommand{\datum}{\today}
\newcommand{\studentA}{Robin Klusman}
\newcommand{\uvanetidA}{10675671}
\newcommand{\studentB}{Maico Timmerman}
\newcommand{\uvanetidB}{10542590}
%\newcommand{\studentC}{Boudewijn Braams}
\newcommand{\uvanetidC}{10401040}
%\newcommand{\studentD}{Govert Verkes}
\newcommand{\uvanetidD}{10211748}
%\newcommand{\studentE}{Naam student 5}
\newcommand{\uvanetidE}{UvAnetID student 5}

% ===================================  ===================================

\begin{document}
\thispagestyle{firststyle}
\begin{center}
    \rule{\linewidth}{0.5pt} \\[0.4cm]
    {\LARGE \bfseries \titel}
    \rule{\linewidth}{0.5pt} \\[0.2cm]
    \textsc{\Large \opdracht}\\[0.2cm]
    \vspace{0.2cm}
    {\large \datum  \\[0.4cm]}

    \begin{minipage}{0.4\textwidth}
        \begin{flushleft}
            \emph{Authors:}\\
            {\studentA \\ {\small \uvanetidA \\[0.2cm]}}
            \ifthenelse{\isundefined{\studentB}}{}{\studentB \\ {\small \uvanetidB \\[0.2cm]}}
            \ifthenelse{\isundefined{\studentC}}{}{\studentC \\ {\small \uvanetidC \\[0.2cm]}}
            \ifthenelse{\isundefined{\studentD}}{}{\studentD \\ {\small \uvanetidD \\[0.2cm]}}
            \ifthenelse{\isundefined{\studentE}}{}{\studentE \\ {\small \uvanetidE \\ [0.2cm]}}
        \end{flushleft}
    \end{minipage}
    ~
    \begin{minipage}{0.4\textwidth}
        \begin{flushright}
            \emph{Supervisor:} \\
            \docent \\[0.2cm]
            \emph{Course:} \\
            \cursus \\[0.2cm]
            \emph{Course code:} \\
            \vakcode \\[0.2cm]
        \end{flushright}
    \end{minipage}\\[1 cm]
\end{center}


% =================================== FRONT PAGE ===================================

\clearpage
\vspace*{0.5cm}
\tableofcontents
\clearpage



% =================================== MAIN TEXT ===================================

\section{Abstract}

A distributed file system is a client and server system where clients can access
data located on the servers. The servers can be spread over multiple locations
and should still provide reliable access to files. To do so there are a few
requirements that need to be met: universality, performance, and fault
tolerance.  In this paper our aim is to examine what trade-offs designers of
distributed file systems made in order to make their file systems work. To do
this we will discuss three example distributed file systems, LOCUS, Andrew and
Ceph and look at how they met these requirements or whether they needed to
compromise in some way. We will also compare the three to see how distributed
file systems developed over the years.

\section{Introduction}

A file system is the part of a computer system that manages long term storage of
data.  The most common file system is a centralised file system. This type of
file system is the one we usually find in a personal computer, here files are
stored on a single storage device at a single location.

Another type of file system is the distributed file system, in these file
systems data can be stored on multiple storage devices that also do not
necessarily need to be at the same location. These file systems are useful when
multiple people need to share their data quickly and securely.

In this paper our aim is to find out what the trade-offs are exactly  when
choosing a distributed file system over a centralised file system. And also how
distributed file systems developed over the years. We will do this by looking at
three examples and comparing the properties of those to each other and to a
centralised file system.

In the next section we take a closer look at what a distributed file system is,
and why it is useful. Then we will discuss in short what the requirements are
for a distributed file system to be effective and useful. In the section
thereafter we will look at a couple of examples of distributed file systems to
see what trade-offs they have had to make. Then the most important trade-offs
will be summarised and discussion in the last section of this paper.

\section{The Distributed File System}

In this section we take a more in depth look at what a distributed file system
exactly is, and why it is useful. A distributed file system is a system of
servers and clients (in some cases devices can be both server and client at the
same time) where data is spread over multiple storage devices (servers) but can
be accessed by clients as though it were on their own device.\cite{concepts}

A distributed file system is particularly convenient in large corporations where
employees each have their own account that they can use to log in to any
computer at the office. The employees will want to access their files regardless
of which computer they are logged in from. A distributed file system can make
sure every employee will have access to their files as though the files were on
the computer currently being used.

\section{Requirements}

Here we will discuss in short what the requirements are of a distributed file
system in order for it to be effective. First of all a user must be able to
access the files they are trying to access in a way that is (preferably) no
different accessing files on a centralised file system. And a user must also be
able to access their files located on the distributed file system from any
machine within the system.\cite{concepts} Another very important requirement is
of course the performance of the system. If a file system is too slow, it
becomes impractical to use since every time a client requests access to a file
they would have to wait for it. Another key concern is the system's fault
tolerance. Being fault tolerant means that the system will not completely crash
or shut down when a fault occurs, instead it should either manage the fault or
continue functioning with reduced functionality until the fault is
solved.\cite{concepts}

\section{LOCUS}
The LOCUS distributed file system was developed in the early 1980s.

\section{Andrew}

The distributed file system used by the Andrew distributed computing environment
(hereafter referred to as simply Andrew) started development in 1983, just after
the LOCUS file system was developed. In creating Andrew the designers main goals
were performance and scalability. This means that Andrew had to provide a
relatively good performance and also make sure the performance does not suffer
too much when scaling the system.\cite{andrew}

Andrew works by using dedicated servers and clients. The clients each have local
storage and are also able to access the shared storage on the servers. The
Servers are collectively responsible for managing the shared
storage.\cite{concepts} When a client wants to retrieve a file from the shared
storage it will send a request to a server. Servers continuously listen for
client connections and as soon as a client wants to connect the server creates a
new process to handle the communication with that client. The process is only
terminated when the client disconnects.\cite{andrew} Servers contain both data
and links to data that exists on other servers.

In Andrew universality is achieved through caching. When a file is requested it
is put into local storage before the user can manipulate the file. This means
that the way a user interacts with file is much the same as on a centralised
file system, since the file is actually copied to local memory. The only
difference is the time it takes to load a not yet cached file.\cite{concepts}

Caching also plays a great part in performance. Because files are cached
locally, data transfers between clients and servers occur less frequently. Since
data transfers to servers are very slow and local storage accesses are
relatively fast this significantly increases performance. In addition, the
performance when scaling the system is also very good since the clients can
execute programs that are stored locally without any server interaction, thus
not overloading the server with requests.\cite{andrew}




\section{Ceph}


\section{Discussion}



% =================================== REFERENCES ===================================

\clearpage
\bibliographystyle{unsrt}
\bibliography{bib}

\end{document}
